\section{Training-source audits and representation diagnostics}
\paragraph{Hugging Face parity.} At revision \texttt{2937df7f} (full revision in the provenance manifest), the EgoNormia parquet contains 1,743 unique IDs, all present in GitHub; 110 GitHub IDs are absent. Shared behaviors, justifications, gold, sensible sets, and taxonomy agree after removing null padding in parquet taxonomy keys. Descriptions differ for 27 shared items. This establishes the exact membership gap, but not why the release omitted those items. The analyses consistently use GitHub text/gold. Two GitHub descriptions are missing. Five label occurrences use four unexpected strings; these are retained in the audit and excluded from the seven-label probes without relabeling them.

\paragraph{NormBank.} The downloaded CSV \citep{ziems2023normbank} has 155,423 rows and nine columns, with no parser-detected missing values. Labels are taboo 68,057; normal 59,507; expected 27,859. Provided splits contain 124,920 train, 15,008 dev, and 15,495 test rows. There are 32,059 exact duplicate rows (123,364 distinct rows), and 1,166 exact setting/behavior/constraints keys carry conflicting labels. No such exact context key crosses the provided splits. Exact duplicates and conflicting labels are different findings; neither establishes broader semantic leakage. We recommend reviewing conflicts and deduplicating within a documented policy before training.

\paragraph{NormLens.} The authors' Git-LFS archive \citep{han2023normlens} contains JSON arrays despite the \texttt{.jsonl} suffix: 934 high-agreement and 1,049 mid-agreement examples. High-agreement labels are inappropriate 187, appropriate 350, impossible 397. Mid-agreement unique-majority labels are 254/307/371 respectively, with 117 tied-majority examples retained as ties. High-agreement examples have unanimous \emph{available} votes, but four have just one retained vote; unanimity does not mean five votes for every example. For repeated normalized action text across multiple images, the high-agreement subset has 69 groups and 170 rows; 25 groups have differing judgments, and 29/170 rows (17.1\%) disagree with a modal text-group label under minimum-disagreement tie handling. Combining both subsets after excluding image-level ties gives 162 groups, 418 rows, and 82/418 deviations (19.6\%). These descriptive associations are not paired text-only versus image-conditioned human judgments or proof of image causality.

\paragraph{Contextual language features.} Frozen \texttt{all-mpnet-base-v2} embeddings use the model's default 384-token limit, normalized outputs, and fold-local scaling plus logistic regression. The full five-choice source-grouped action probe scores 54.40\% (1,008/1,853). Seven-label macro AUROC is 0.773 for nonempty candidate options ($n=7{,}412$) and 0.569 for descriptions predicting correct-option labels ($n=1{,}851$). These tasks have different targets and units; their gap does not isolate a difference in modality quality. The t-SNE plot uses only single-label options and is exploratory, not a separability test.

\paragraph{Visual features with matched targets.} We select 400 shared-release IDs by sorting SHA-256 hashes of a fixed prefix plus ID, then extract frozen CLIP ViT-B/32 features. Each downloaded horizontal strip is divided into five equal-width tiles. Each tile is padded, preserving aspect ratio and its full contents, to $224\times224$; we average the five unnormalized feature vectors and normalize the resulting vector. This avoids silently center-cropping away most of a strip, but destroys temporal order and compresses fine interaction details. On identical examples, folds, and correct-option targets, macro AUROC is 0.561 for CLIP, 0.626 for TF-IDF+LSA descriptions, and 0.552 for MPNet descriptions. Fold-mean macro AUROCs are 0.560/0.635/0.552. The TF-IDF+LSA baseline uses 128 fold-local SVD components. All are fixed-hyperparameter diagnostics, not optimized capacity comparisons. Privacy has only 12 positives in this sample, so its high CLIP AUROC should not be treated as stable evidence of general superiority.

\paragraph{Implications and limits.} Generic pooled visual features do not establish the curriculum's desired social reasoning capability. The results motivate preserving interaction details, task-native labels, and label uncertainty. VideoNorms' linked public repository contains only a README at the inspected commit; no new VideoNorms annotation analysis was possible. The team reports that it has contacted the authors to request access and expects to receive the annotations soon; receipt and a delivery date are not yet confirmed. One participant completed all 20 packet items in two sequential blocks; action agreement was 8/20, 8/20, and 9/20. Feedback on the first block preceded the second. The main section reports this exploratory exercise; its full protocol, responses, and limitations are supplied separately.
